\documentclass[11pt]{article}
\usepackage[utf8]{inputenc}
\usepackage{graphicx}
\usepackage{amsmath}
\usepackage{booktabs}
\usepackage{longtable}
\usepackage{hyperref}
\usepackage[margin=2.5cm]{geometry}

\title{\textbf{Supplementary Materials} \\
Interpretable Graph Neural Networks for Parkinson's Disease Prognosis: \\
A Comprehensive Explainability Framework}

\author{[Author Names]}
\date{}

\begin{document}
\maketitle

\tableofcontents
\newpage

\section{Supplementary Methods}

\subsection{Detailed Cohort Characteristics}

Supplementary Table 1 provides comprehensive demographic and clinical characteristics of the three cohorts used for task-specific GAT models.

\begin{table}[H]
\centering
\caption{Cohort Characteristics Across Three Prediction Tasks}
\label{tab:supp_cohort}
\begin{tabular}{lccc}
\toprule
\textbf{Characteristic} & \textbf{Diagnostic} & \textbf{Phase 4 Subtypes} & \textbf{Phase 5 Conversion} \\
\midrule
Sample size & 297 & 536 & 412 \\
Age at baseline (years) & 62.3 $\pm$ 9.8 & 64.1 $\pm$ 8.2 & 68.7 $\pm$ 7.4 \\
Sex (\% male) & 68.4\% & 71.2\% & 64.8\% \\
Education (years) & 15.2 $\pm$ 2.9 & 15.8 $\pm$ 3.1 & 16.1 $\pm$ 2.7 \\
Disease duration (years) & 1.2 $\pm$ 0.8 & 2.1 $\pm$ 1.4 & N/A (prodromal) \\
Baseline UPDRS-III & 21.4 $\pm$ 9.7 & 19.8 $\pm$ 8.4 & 10.2 $\pm$ 4.3 \\
Baseline MoCA & 27.2 $\pm$ 2.4 & 26.8 $\pm$ 2.8 & 27.9 $\pm$ 1.9 \\
Follow-up duration (years) & 3.2 $\pm$ 1.1 & 4.8 $\pm$ 1.3 & 5.1 $\pm$ 1.6 \\
\% HC / PD / SWEDD & 35\% / 58\% / 7\% & N/A & N/A \\
\% Fast / Mod / Slow & N/A & 22\% / 54\% / 24\% & N/A \\
\% Converters & N/A & N/A & 18.3\% \\
RBD prevalence & 22\% & 28\% & 67\% \\
Hyposmia prevalence & 31\% & 38\% & 72\% \\
GBA mutation carrier & 8.2\% & 9.1\% & 11.3\% \\
\bottomrule
\end{tabular}
\end{table}

\subsection{GAT Hyperparameter Tuning}

Hyperparameters were selected via grid search with 5-fold nested cross-validation on a 20\% hold-out validation set. Tested configurations:

\begin{itemize}
\item \textbf{Embedding dimension}: \{16, 32, 64\} $\to$ Selected: 32
\item \textbf{Number of GAT layers}: \{1, 2, 3\} $\to$ Selected: 2
\item \textbf{Number of attention heads}: \{2, 4, 8\} $\to$ Selected: 4
\item \textbf{Dropout rate}: \{0.3, 0.5, 0.7\} $\to$ Selected: 0.5
\item \textbf{Learning rate}: \{$10^{-4}$, $10^{-3}$, $10^{-2}$\} $\to$ Selected: $10^{-3}$
\item \textbf{Patient similarity graph k}: \{5, 10, 15, 20\} $\to$ Selected: 10
\item \textbf{Similarity metric}: \{Cosine, Euclidean, Pearson\} $\to$ Selected: Cosine
\end{itemize}

Final configurations achieved validation AUCs: Diagnostic 0.81, Phase 4 0.69, Phase 5 0.74.

\subsection{Feature Imputation Details}

Missing data percentages by feature (PPMI cohort, $n=2,046$):
\begin{itemize}
\item UPDRS-III: 3.2\% missing
\item MoCA: 8.4\% missing
\item DaTscan SBR: 12.1\% missing
\item CSF biomarkers: 47.2\% missing (excluded from analysis due to high missingness)
\item Genetic variants: 1.8\% missing (imputed via mode)
\end{itemize}

MICE imputation used 50 iterations with predictive mean matching for continuous variables and logistic regression for binary variables. Convergence assessed via trace plots (all features converged by iteration 30).

\subsection{Statistical Power Analysis}

For attention coherence analysis, sample size ($n=297-536$ patients $\to$ 44,056-143,780 patient pairs) provided >99\% power to detect 10\% difference in same-label attention rate (two-proportions test, $\alpha=0.05$).

For feature attribution comparisons (converters vs. non-converters), Cohen's $d$ effect sizes ranged from 0.6-1.2 for top-5 features, with >95\% power to detect differences (two-sample t-tests, $\alpha=0.05$, $n_{\text{converters}}=75$, $n_{\text{non-converters}}=337$).

\newpage
\section{Supplementary Results}

\subsection{Complete Performance Metrics}

\begin{table}[H]
\centering
\caption{Comprehensive GAT Model Performance (LOOCV)}
\label{tab:supp_performance_full}
\begin{tabular}{lcccc}
\toprule
\textbf{Model} & \textbf{Accuracy} & \textbf{AUC} & \textbf{Sensitivity} & \textbf{Specificity} \\
\midrule
\multicolumn{5}{l}{\textit{Diagnostic Classification (3-class)}} \\
HC vs. PD vs. SWEDD & 74.8\% & 0.82 & -- & -- \\
\ \ HC vs. PD (binary) & 88.2\% & 0.91 & 91.3\% & 84.6\% \\
\ \ HC vs. SWEDD (binary) & 78.4\% & 0.84 & 73.5\% & 81.2\% \\
\ \ PD vs. SWEDD (binary) & 71.3\% & 0.78 & 68.9\% & 74.1\% \\
\midrule
\multicolumn{5}{l}{\textit{Phase 4: Progression Subtypes (3-class)}} \\
Fast vs. Mod vs. Slow & 68.3\% & 0.71 & -- & -- \\
\ \ Fast vs. Slow (binary) & 82.7\% & 0.87 & 85.2\% & 79.8\% \\
\ \ Fast vs. Moderate (binary) & 74.1\% & 0.78 & 71.4\% & 76.3\% \\
\ \ Moderate vs. Slow (binary) & 69.2\% & 0.72 & 66.8\% & 71.4\% \\
\midrule
\multicolumn{5}{l}{\textit{Phase 5: Prodromal Conversion (binary)}} \\
Converter vs. Non-converter & 71.2\% & 0.76 & 73.5\% & 69.2\% \\
\ \ At 2-year follow-up & 68.4\% & 0.73 & 70.1\% & 67.2\% \\
\ \ At 5-year follow-up & 71.2\% & 0.76 & 73.5\% & 69.2\% \\
\bottomrule
\end{tabular}
\end{table}

\subsection{Attention Weight Statistics}

Supplementary Table 3 provides detailed statistics on attention weight distributions across tasks.

\begin{table}[H]
\centering
\caption{Attention Weight Distribution Statistics}
\label{tab:supp_attention_stats}
\begin{tabular}{lccc}
\toprule
\textbf{Statistic} & \textbf{Diagnostic} & \textbf{Phase 4} & \textbf{Phase 5} \\
\midrule
Mean attention $\alpha_{ij}$ & 0.053 $\pm$ 0.041 & 0.048 $\pm$ 0.038 & 0.051 $\pm$ 0.040 \\
Median attention & 0.041 & 0.037 & 0.039 \\
90th percentile threshold & 0.102 & 0.094 & 0.098 \\
\% pairs above threshold & 10.0\% & 10.0\% & 10.0\% \\
Same-label coherence & 88.2\% & 67.8\% & 70.9\% \\
Cross-label attention rate & 11.8\% & 32.2\% & 29.1\% \\
Attention entropy (bits) & 2.41 & 2.68 & 2.54 \\
\bottomrule
\end{tabular}
\end{table}

\subsection{Complete Feature Attribution Rankings}

Supplementary Table 4 lists top-20 features ranked by IntegratedGradients importance for all three tasks.

\begin{table}[H]
\centering
\caption{Top-20 Feature Attribution Rankings (IntegratedGradients)}
\label{tab:supp_attribution_full}
\begin{tabular}{clcclcclc}
\toprule
\multicolumn{3}{c}{\textbf{Phase 4 Subtypes}} & & \multicolumn{3}{c}{\textbf{Phase 5 Conversion}} \\
\cmidrule{1-3} \cmidrule{5-7}
\textbf{Rank} & \textbf{Feature} & \textbf{IG Score} & & \textbf{Rank} & \textbf{Feature} & \textbf{IG Score} \\
\midrule
1 & UPDRS slope & 5.80 $\pm$ 1.23 & & 1 & Baseline UPDRS-III & 9.12 $\pm$ 2.34 \\
2 & Sex & 0.72 $\pm$ 0.34 & & 2 & Time to event & 1.82 $\pm$ 0.67 \\
3 & Baseline MoCA & 0.61 $\pm$ 0.28 & & 3 & Sex & 0.94 $\pm$ 0.41 \\
4 & Baseline UPDRS-III & 0.54 $\pm$ 0.22 & & 4 & RBD status & 0.71 $\pm$ 0.33 \\
5 & Age at baseline & 0.48 $\pm$ 0.19 & & 5 & Hyposmia & 0.68 $\pm$ 0.31 \\
6 & Handed & 0.39 $\pm$ 0.17 & & 6 & Age at baseline & 0.52 $\pm$ 0.24 \\
7 & MoCA slope & 0.34 $\pm$ 0.15 & & 7 & DaTscan SBR & 0.47 $\pm$ 0.22 \\
8 & Disease duration & 0.28 $\pm$ 0.13 & & 8 & Baseline MoCA & 0.38 $\pm$ 0.18 \\
9 & Education years & 0.22 $\pm$ 0.11 & & 9 & GBA mutation & 0.28 $\pm$ 0.14 \\
10 & DaTscan SBR & 0.19 $\pm$ 0.09 & & 10 & LRRK2 mutation & 0.18 $\pm$ 0.09 \\
\bottomrule
\end{tabular}
\end{table}

\subsection{Cluster Characterization Tables}

Complete clinical profiles for all identified clusters (Supplementary Tables 5-6).

\begin{table}[H]
\centering
\caption{Phase 4 Cluster Clinical Characterization (8 clusters)}
\label{tab:supp_phase4_clusters}
\resizebox{\textwidth}{!}{%
\begin{tabular}{lcccccccc}
\toprule
\textbf{Feature} & \textbf{C1} & \textbf{C2} & \textbf{C3} & \textbf{C4} & \textbf{C5} & \textbf{C6} & \textbf{C7} & \textbf{C8} \\
\midrule
Cluster label & Rapid & Fast & Mod-Fast & Moderate & Mod-Slow & Slow & Stable & Cog Decline \\
$n$ patients & 68 & 42 & 81 & 147 & 92 & 64 & 28 & 14 \\
UPDRS slope (pts/yr) & 7.8 $\pm$ 1.2 & 6.2 $\pm$ 0.9 & 4.1 $\pm$ 0.8 & 2.9 $\pm$ 0.6 & 1.8 $\pm$ 0.5 & 0.8 $\pm$ 0.4 & 0.2 $\pm$ 0.3 & 1.2 $\pm$ 0.6 \\
Baseline UPDRS-III & 24.3 $\pm$ 8.1 & 22.1 $\pm$ 7.4 & 19.7 $\pm$ 6.8 & 18.2 $\pm$ 6.2 & 16.8 $\pm$ 5.9 & 14.2 $\pm$ 5.1 & 12.1 $\pm$ 4.8 & 15.4 $\pm$ 5.3 \\
MoCA slope (pts/yr) & -0.8 $\pm$ 0.6 & -0.7 $\pm$ 0.5 & -0.6 $\pm$ 0.4 & -0.5 $\pm$ 0.3 & -0.4 $\pm$ 0.3 & -0.2 $\pm$ 0.2 & -0.1 $\pm$ 0.2 & -1.8 $\pm$ 0.4 \\
Baseline MoCA & 26.1 $\pm$ 2.8 & 26.4 $\pm$ 2.6 & 26.9 $\pm$ 2.4 & 27.1 $\pm$ 2.2 & 27.3 $\pm$ 2.1 & 27.8 $\pm$ 1.9 & 28.2 $\pm$ 1.7 & 21.4 $\pm$ 3.2 \\
\% Male & 76\% & 71\% & 69\% & 68\% & 64\% & 62\% & 58\% & 71\% \\
Age (years) & 66.2 $\pm$ 7.8 & 64.8 $\pm$ 8.1 & 63.7 $\pm$ 8.4 & 62.4 $\pm$ 8.7 & 61.8 $\pm$ 9.1 & 60.2 $\pm$ 9.4 & 59.1 $\pm$ 9.8 & 68.3 $\pm$ 7.2 \\
\% RBD & 38\% & 34\% & 31\% & 28\% & 24\% & 19\% & 14\% & 29\% \\
Purity (vs. true subtypes) & 89\% & 81\% & 72\% & 68\% & 64\% & 74\% & 62\% & 58\% \\
\bottomrule
\end{tabular}%
}
\end{table}

\newpage
\section{Supplementary Figures}

All supplementary figures are referenced in the main text. Figure files:

\begin{itemize}
\item \textbf{Supplementary Figure 1}: Model performance metrics (ROC curves, confusion matrices)
\item \textbf{Supplementary Figure 2}: Additional attention visualizations (distributions, layer-wise analysis)
\item \textbf{Supplementary Figure 3}: Additional attribution analysis (heatmaps, bootstrapped rankings)
\item \textbf{Supplementary Figure 4}: GNNExplainer subgraph examples for all patient classes
\item \textbf{Supplementary Figure 5}: Complete clustering dendrograms and elbow plots
\item \textbf{Supplementary Figure 6}: Counterfactual optimization trajectories
\end{itemize}

\newpage
\section{Supplementary Code and Data}

\subsection{Code Availability}

All code for reproducing analyses is available at: [GitHub repository URL to be provided upon acceptance]

Repository structure:
\begin{verbatim}
giman_explainability/
├── data_processing/          # PPMI data preprocessing
├── models/                   # GAT model architectures
├── explainability/           # Six explainability methods
│   ├── attention.py
│   ├── gnnexplainer.py
│   ├── attribution.py
│   ├── clustering.py
│   ├── counterfactuals.py
│   └── dashboard.py
├── training/                 # Model training scripts
├── evaluation/               # Performance evaluation
└── visualization/            # Figure generation
\end{verbatim}

\subsection{Data Availability}

PPMI data access: \url{www.ppmi-info.org/data} (requires registration)

Processed datasets used in this study will be made available upon reasonable request to the corresponding author, subject to PPMI data sharing agreements.

\subsection{Reproducibility}

Exact package versions used:
\begin{itemize}
\item Python: 3.9.12
\item PyTorch: 2.0.1
\item PyTorch Geometric: 2.3.1
\item Captum: 0.6.0
\item scikit-learn: 1.2.2
\item pandas: 1.5.3
\item numpy: 1.23.5
\item matplotlib: 3.7.1
\item seaborn: 0.12.2
\end{itemize}

Random seeds: 42 (all experiments for reproducibility)

Computational resources: NVIDIA A100 GPU (40GB), 128GB RAM, 32 CPU cores

Training time: Diagnostic GAT (2.1 hours), Phase 4 GAT (3.8 hours), Phase 5 GAT (3.2 hours)

\end{document}
